\begin{table}[t]
\caption{Dynamic power (mW) at matched operand sparsity}
\label{tab:power}
\centering\footnotesize
\setlength{\tabcolsep}{4.5pt}
\begin{tabular}{lcccc}
\toprule
Design & \multicolumn{4}{c}{operand sparsity $s$ (\%)}\\
\cmidrule(l){2-5}
 & 0 & 50 & 75 & 90\\
\midrule
Baseline (2-stage, async reset) & 70 & 65 & 58 & 48\\
\quad + split-enable zero-skip & 72 & 56 & 43 & 29\\
\quad\emph{change vs.\ its control} & \emph{+3\%} & \emph{-14\%} & \emph{-26\%} & \emph{-40\%}\\
\midrule
DSP-packed, gating off & 46 & 45 & 45 & 43\\
\quad + split-enable zero-skip & 47 & 45 & 41 & 37\\
\quad\emph{change vs.\ its control} & \emph{+2\%} & \emph{-0\%} & \emph{-9\%} & \emph{-14\%}\\
\midrule
clock component (baseline) & 20 & 20 & 20 & 20\\
\bottomrule
\end{tabular}
\\[2pt]
\raggedright\scriptsize Vivado \texttt{report\_power}; confidence Medium.
Each design is compared against its own structural control at the same
sparsity, because the ungated baseline itself loses 31\,\% of its power
between $s=0$ and $s=0.9$ from reduced multiplier switching alone.
The DSP-packed rows understate the gating benefit: see
Section~\ref{sec:results}-E.
\end{table}
